\scanpage{337}\mspage{253}
% l. 1–4
$(\alpha)$ en composant avec $F(j')\unsure{\mathbin{-}|\longrightarrow}X(j')$ soit commutatif. On \deleted{choisit} fait ce choix pour tt objet $j'$ de $J$. Appliquons alors $\mathrm F'1$ \inserted{à $F$ et} au système de flèches
% l. 5–7
$F(j)\xrightarrow{v(j)}X(j)$, on trouve $G$ et un système de flèches \deleted{$\unsure{w(j)|v(j)}$}
\[
 w(j):X(j)\unsure{\mathbin{-}|\longrightarrow}G(j)
\]
% l. 8–11
de telle façon que \deleted{pour tte flèche} les flèches composées
\diagram{\begin{tikzcd}[ampersand replacement=\&,column sep=small,row sep=small]
 w(j)v(j):\&F(j)\arrow[rr]\arrow[dr,dashed]\& \&G(j)\\
 \& \&X(j)\arrow[ur,dashed]\&
\end{tikzcd}}
définissent un hom de sections\unsure{,|{}}. Je dis que cet hom. \unsure{à|a} la propriété voulue pour $\mathrm F3_{\unsure{\pi|\kappa}}$, \unsure{i.e|i.e.|vz} que pour tt $\alpha\in A$, le système des composés
% l. 12–16
\[
 F_{\alpha}(j)\unsure{\mathbin{-}|\longrightarrow} F(j)\unsure{\mathbin{-}|\longrightarrow}G(j)
\]
(\deleted{\unsure{pour}} $j\in\operatorname{Ob}J$ variable) définit un hom de sections $F_{\alpha}\longrightarrow G$. En effet, \deleted{\unsure{évident}} \inserted{il suffit} pour cela de \unsure{considérer}, pour une flèche $\lambda$ de $J$, le
% archive sheet number: 336
